\documentclass[11pt]{article}

%==============================================================================
%  LPS Local Dimension Regularization -- Implementation Notes
%  Test and experiment specifications + a VALENCIA-inspired non-manifold set.
%  Companion to dev/notes/foundations/lps_local_dimension_regularization.tex (the expository
%  note) and lps_experimental_plan_2026-06-09.tex (the experimental plan).
%  All figures are original schematics; no data or figures are reproduced.
%==============================================================================

\usepackage[margin=1in]{geometry}
\usepackage{amsmath,amssymb,amsthm,mathtools}
\usepackage{graphicx}
\usepackage{xcolor}
\usepackage{booktabs}
\usepackage{enumitem}
\usepackage{microtype}
\usepackage{caption}
\usepackage[most]{tcolorbox}
\usepackage{datetime2}
\DTMsetup{showzone=false}        % timestamp shown in US Eastern time (build sets TZ); see latexmkrc
\usepackage[hidelinks]{hyperref}
\usepackage[section]{placeins}
\widowpenalty=10000
\clubpenalty=10000
\raggedbottom

\definecolor{ink}{HTML}{24303D}
\definecolor{accent}{HTML}{2563EB}
\definecolor{accent2}{HTML}{B45309}
\definecolor{accent3}{HTML}{15803D}
\captionsetup{font=small,labelfont=bf}
\setlist{topsep=3pt,itemsep=2pt,parsep=1pt}

\theoremstyle{definition}
\newtheorem{definition}{Definition}[section]
\newtheorem{dgp}{DGP}[section]

\newtcolorbox{plain}{
  enhanced, colback=accent!4, colframe=accent!55!black, boxrule=0.4pt, arc=2pt,
  left=7pt, right=7pt, top=4pt, bottom=4pt, title=In plain terms,
  fonttitle=\bfseries\sffamily\small, coltitle=accent!50!black,
  attach title to upper={\quad}}

\newtcolorbox{decision}{
  enhanced, colback=accent3!5, colframe=accent3!55!black, boxrule=0.4pt, arc=2pt,
  left=7pt, right=7pt, top=4pt, bottom=4pt, title=Decision rule,
  fonttitle=\bfseries\sffamily\small, coltitle=accent3!45!black,
  attach title to upper={\quad}}

\newcommand{\R}{\mathbb{R}}
\newcommand{\E}{\mathbb{E}}
\newcommand{\supp}{\operatorname{supp}}
\newcommand{\argmin}{\operatorname*{arg\,min}}
\newcommand{\tr}{\operatorname{tr}}

% test/experiment field labels
\newcommand{\field}[1]{\smallskip\noindent\textbf{\textsf{#1.}}\ }

\newcommand{\fig}[3]{%
  \begin{figure}[htbp]\centering
  \includegraphics[width=#2]{../../../notes/foundations/figures_ldr/#1}%
  \caption{#3}\label{fig:#1}\end{figure}}

\title{\textbf{LPS Local Dimension Regularization: Implementation Notes}\\[4pt]
\large Precise test and experiment specifications, and a VALENCIA-inspired
non-manifold example set with known ground-truth dimension}
\author{\small Internal implementation note --- LPS / \texttt{geosmooth} program\\[2pt]
\footnotesize source: \nolinkurl{~/current_projects/geosmooth/dev/methods/lps/specs/lps_local_dimension_regularization_implementation_notes.tex}}
\date{\DTMnow\ \textnormal{\footnotesize ET}}

\begin{document}
\maketitle

\begin{abstract}\noindent
This is the implementation companion to the expository note
\emph{Local Dimension Regularization in Local Polynomial Smoothing} and to the
\emph{LPS Experimental Plan}. It specifies, precisely enough to implement without
further design, two distinct kinds of evaluation. \textbf{Tests} are
deterministic, pass/fail checks of the local-dimension-field machinery, run on
\emph{synthetic} data whose ground-truth dimension field we designed and therefore
know exactly. \textbf{Experiments} are exploratory studies of the \emph{real}
VALENCIA 16S data (the France et al.\ collection, $n\approx 13{,}231$), restricted
to relative abundances on the simplex; they have no ground truth, and their
purpose is to characterize the data and to \emph{guide next steps}. The two are
connected by a single loop: the experiments measure the structure of the real
data --- which support patterns (naive strata) occur, and how the mass
concentrates near lower-dimensional faces --- and those measurements parameterize
a family of \emph{purely synthetic} generators that reproduce that structure while
carrying a known local-dimension field. Those generators become the expanded
non-manifold example set on which the tests run. Every test and experiment is
stated with its objective, exact procedure or configuration, the statistic
computed, and a numeric acceptance criterion or decision rule.
\par\medskip\noindent\footnotesize\textit{On the figures.} All figures are
original schematics drawn for this note; the ``VALENCIA'' panels are illustrative
sketches of the \emph{shape} of the outputs the experiments will produce, not
computed results (the data are not used in producing this document).
\end{abstract}

{\small\tableofcontents}

%==============================================================================
\section{Scope, and the tests--experiments distinction}\label{sec:scope}
%==============================================================================

The expository note argued, conceptually, that the local chart dimension in LPS
should be estimated as a \emph{field} over the neighbor graph and regularized with
an $\ell_1$ (total-variation) penalty so that it denoises within strata while
keeping the jumps between strata sharp. This note turns that argument into an
implementable validation program. The single most important design decision is
to keep two kinds of evaluation rigorously separate.

\begin{itemize}[leftmargin=1.4em]
\item \textbf{Tests} (Section~\ref{sec:tests}) are deterministic and pass/fail.
They run on synthetic data whose true local-dimension field is \emph{known by
construction}, so ``correct'' is well defined and a continuous-integration suite
can enforce it. They answer: \emph{does the estimator do what it claims?}
\item \textbf{Experiments} (Section~\ref{sec:experiments}) are exploratory. They
run on the real VALENCIA data, which has \emph{no} ground-truth dimension, so they
produce descriptive statistics and figures rather than verdicts. They answer:
\emph{what does the real data look like, and what should we build and try next?}
\end{itemize}

These must not be conflated. A statistic measured on real data (an estimated
effective dimension, say) is evidence about the data and about our estimator's
behavior, but it can never be a pass/fail test, because there is nothing to check
it against. Conversely, a test that passed on a synthetic toy says nothing about
realism. The two are joined by the loop in Figure~\ref{fig:fig_flow.pdf}: the
experiments characterize the real data; that characterization parameterizes
synthetic generators (Section~\ref{sec:synthetic}) that carry a known dimension
field; and the tests run on those generators. Real data supplies \emph{realism};
synthetic data supplies \emph{ground truth}; neither alone suffices.

\fig{fig_flow.pdf}{\linewidth}{How the pieces fit. The real VALENCIA data is
explored (experiments E1--E3, no ground truth) to characterize its stratified
structure. That structure parameterizes synthetic generators with a \emph{known}
local-dimension field, which become the example set on which the deterministic
tests T1--T9 run. Test outcomes and experiment findings together feed the
next-step decisions.}

\begin{plain}
We are checking a tool that guesses the ``local dimension'' of a data cloud.
To check a tool you need an answer key, and real biological data does not come
with one. So we do two things. We \emph{explore} the real microbiome data to learn
its shape (no answer key --- this guides what to build). Then we \emph{build}
fake data that has the same shape but where we wrote the answer key ourselves, and
we grade the tool against that. The fake data is realistic because the exploration
told us how; it is gradeable because we made it.
\end{plain}

%==============================================================================
\section{The object under test, and shared conventions}\label{sec:object}
%==============================================================================

We recall the estimator precisely, because the tests are written against it. Over
the neighbor graph $G=(V,E)$ on $n$ anchors, the local-dimension field
$d=(d_1,\dots,d_n)\in\{1,\dots,D_{\max}\}^n$ minimizes
\begin{equation}\label{eq:field}
  \widehat d \;=\; \argmin_{d}\;
  \sum_{i\in V} c_i(d_i)\;+\;\lambda\!\!\sum_{(i,j)\in E}\!\! w_{ij}\,\rho(d_i,d_j),
\end{equation}
where the per-anchor cost $c_i(d)$ is a complexity-penalized fit of dimension $d$
to anchor $i$'s neighborhood (e.g.\ the residual eigen-energy beyond $d$ principal
components measured at the smoother's bandwidth scale, plus a penalty), $w_{ij}$
are graph weights, $\lambda\ge 0$ is the regularization strength, and the pairwise
penalty is the total variation $\rho(d_i,d_j)=|d_i-d_j|$ (the $\ell_2$ comparison
penalty $(d_i-d_j)^2$ appears only in tests that contrast the two). Because the
labels are linearly ordered integers and $\rho$ is convex, \eqref{eq:field} is
minimized \emph{exactly} by a single graph cut (Ishikawa). The \emph{soft face} at
scale $h$ folds near-zero coordinates into a lower stratum, and the effective
dimension is capped by the support: $d_i\le |A_i|-1$ where $A_i=\supp(x_i)$.

\field{Conventions} The tests inherit the conventions of the experimental plan and
restate only what is specific here.
\begin{itemize}[leftmargin=1.4em,nosep]
\item \emph{RNG.} Every synthetic dataset is generated after an explicit seed;
seeds and the package commit are stored with each result. Tests are bitwise
reproducible.
\item \emph{Ground truth.} For a synthetic generator the true field is the pair
$(A_i,k_i)$ at each sample: the support pattern $A_i$ (the naive stratum) and the
designed effective dimension $k_i$. The ``true local dimension'' at scale $h$ is
$k_i$ when $h$ exceeds the tube thickness $\rho$ and $|A_i|-1$ below it
(Section~\ref{sec:synthetic}); tests fix $h$ in the regime they intend.
\item \emph{Metrics.} Dimension misclassification rate $\frac1n\sum_i\mathbf 1\{\widehat
d_i\neq k_i\}$; boundary blur (defined in T2/T4); and, for the downstream test,
LPS probability/continuous Truth-RMSE.
\item \emph{Coordinates.} Unless a test is specifically about the coordinate
choice (T6), dimension is estimated in the cube/raw-abundance coordinate $\xi$, not
the log-ratio $u$ (the expository note's coordinate trap).
\end{itemize}

\begin{plain}
Equation~\eqref{eq:field} is the whole tool in one line: each anchor proposes a
dimension from its own data ($c_i$), neighbors are pulled toward agreement
($\rho$), and because dimension is a small whole number the best compromise can be
found exactly. The tests below poke each part of that sentence.
\end{plain}

%==============================================================================
\section{Part A --- Test specifications (synthetic, deterministic, pass/fail)}
\label{sec:tests}
%==============================================================================

Each test follows the experimental-plan template: a falsifiable \textsf{Claim};
the \textsf{DGP} (a named synthetic generator from Section~\ref{sec:synthetic},
whose true field $(A_i,k_i)$ is known); the \textsf{Configuration} (every pinned
argument); the \textsf{Statistic}; a numeric \textsf{Acceptance} criterion; and
\textsf{Safeguards} against a confounded result. The DGP names (stratum mixture,
tube, seam, cone) are defined precisely in Section~\ref{sec:synthetic}; here we
only name the regime each test needs.

\subsection{T1 --- the field solver returns the exact global optimum}
\field{Claim} For the total-variation penalty, the graph-cut solver of
\eqref{eq:field} returns the exact global minimum.
\field{DGP} Not a data generator: random small instances. Graphs with $n\le 10$
nodes, random edges, random integer costs $c_i(d)$ for $d\in\{1,2,3\}$, random
$\lambda$, over $R=200$ seeds.
\field{Configuration} Solve by the graph-cut path; independently minimize by
exhaustive enumeration of all $3^{n}$ labelings.
\field{Statistic} Energy gap $\Delta = E_{\text{cut}} - E_{\text{brute}}$.
\field{Acceptance} $\Delta = 0$ exactly on every instance (the cut is the global
optimum, not an approximation).
\field{Safeguards} Include instances with cost ties and with $\lambda$ large
enough to force a single constant label and small enough to recover the unary
argmin, so both extremes are exercised.

\subsection{T2 --- $\ell_1$ keeps the boundary sharp; $\ell_2$ blurs it}
\field{Claim} At a true dimension jump, the TV field transitions over a narrow
band of anchors while the Laplacian ($\ell_2$) field ramps over a wide one.
\field{DGP} The \emph{seam} fragment: two strata of dimensions $k_1\neq k_2$
meeting along a known boundary, with i.i.d.\ per-anchor estimation noise.
\field{Configuration} Estimate the field with $\rho=|\cdot|$ (TV) and with
$\rho=(\cdot)^2$ ($\ell_2$), at $\lambda$ values \emph{calibrated to equal
interior denoising} (matched misclassification far from the seam), so the
comparison is only about the boundary.
\field{Statistic} Boundary width $B$: the number of anchors whose fitted dimension
lies strictly between $k_1$ and $k_2$ (TV labels are integers, so for TV this is
the count of anchors not yet snapped to either level within a fixed corridor of
the seam).
\field{Acceptance} $B_{\ell_1}\le 2$ and $B_{\ell_2}\ge 5\,B_{\ell_1}$.
\field{Safeguards} Calibrate the two penalties to equal interior denoising first;
identical noise realization for both arms (paired).

\subsection{T3 --- regularization reduces dimension misclassification}
\field{Claim} The TV field has lower dimension-misclassification than independent
per-anchor estimation.
\field{DGP} The stratum-mixture generator in a near-homogeneous regime (few
strata, well-separated $k_A$), small neighborhoods (high per-anchor variance).
\field{Configuration} Independent argmin of $c_i$ vs.\ TV field; $\lambda$ fixed
(and, separately, $\lambda$ chosen by T9's held-out rule).
\field{Statistic} Misclassification rate $\frac1n\sum_i\mathbf 1\{\widehat d_i\neq
k_i\}$, averaged over $R\ge 30$ replicates.
\field{Acceptance} TV misclassification is below independent by a predeclared
relative margin (e.g.\ $\ge 30\%$ at the smallest neighborhood size).
\field{Safeguards} Same data for both arms; report the full misclassification
distribution, not only the mean.

\subsection{T4 --- the boundary safeguard (regularization must not erase a real jump)}
\field{Claim} TV regularization denoises without smearing a true dimension jump;
$\ell_2$ regularization fails this and is the negative control.
\field{DGP} The seam fragment with a known boundary (as T2).
\field{Configuration} Independent, TV, and $\ell_2$ fields at matched interior
denoising.
\field{Statistic} Boundary-straddling misclassification: the error rate among
anchors within graph radius $r$ of the seam.
\field{Acceptance} TV increases boundary-straddling misclassification over the
independent estimate by $\le 5$ percentage points, \emph{and} $\ell_2$ exceeds
that bound (demonstrating the discriminator works).
\field{Safeguards} This test is mandatory and is the one a naive smoother fails;
if $\ell_2$ does not fail it, the seam is too weak --- strengthen the jump.

\subsection{T5 --- scale-dependence and persistence}
\field{Claim} The effective dimension equals the designed $k_A$ at scales above
the tube thickness $\rho$ and $|A|-1$ below it; the bandwidth-tied estimate
returns $k_A$.
\field{DGP} The tube fragment with known $\rho$ and $k$.
\field{Configuration} Compute $c_i(d;h)$ over a geometric band of scales $h$;
report the dimension-vs-scale profile and the persistence (range of $h$ over which
the chosen dimension is stable).
\field{Statistic} Recovered dimension-vs-scale curve compared to the designed step
at $\rho$; the value at $h=h_{\text{band}}$.
\field{Acceptance} The recovered step location is within a tolerance of $\rho$,
the high-scale plateau equals $k_A$, and the bandwidth-tied estimate equals $k_A$.
\field{Safeguards} $\rho$ and $k$ known; the test is a curve match, not a single
number.

\subsection{T6 --- estimate dimension in $\xi$, not in $u$ (the coordinate trap)}
\field{Claim} Local PCA in the cube coordinate $\xi$ recovers the tube dimension
$k$; local PCA in the log-ratio coordinate $u$ inflates it near a face.
\field{DGP} A face-approaching tube: a stratum whose mass concentrates toward a
lower face (some present coordinates near the detection limit), known $k$.
\field{Configuration} Estimate the effective rank by local PCA in $\xi$ and,
separately, in $u=\log\xi$.
\field{Statistic} Estimated effective rank in each coordinate.
\field{Acceptance} The $\xi$-estimate equals $k$ within tolerance; the
$u$-estimate is demonstrably larger ($\ge k+1$), confirming the inflation.
\field{Safeguards} Known $k$; concentration deliberately placed near a face so the
two coordinates disagree.

\subsection{T7 --- the support cap and the soft face}
\field{Claim} (a) The estimated dimension never exceeds the support cap; (b) the
soft-face active set is monotone in scale and reduces to the hard support as
$h\to 0$.
\field{DGP} The stratum-mixture generator with structural zeros \emph{and}
near-zeros (taxa at the detection limit).
\field{Configuration} Vary the soft-face scale $h$ / threshold $\tau_h$.
\field{Statistic} $\max_i\big(\widehat d_i-(|A_i|-1)\big)$; and, for each sample,
the effective active set $A_h(x)$ as a function of $h$.
\field{Acceptance} The cap is never violated ($\max\le 0$); $A_h$ is nested
(monotone) in $h$; $A_{h\to 0}=\supp(x)$.
\field{Safeguards} Include samples sitting exactly on faces (cap tight) and
samples in tubes near faces (cap loose) so both regimes are tested.

\subsection{T8 --- the frontier (no impossible high-dimensional islands)}
\field{Claim} The field contains no geometrically impossible isolated
high-dimensional anchor (one whose neighbors are all lower-dimensional), which a
proper stratification forbids by lower-semicontinuity.
\field{DGP} The cone fragment: a stratum whose dimension drops at a known
singular point/locus (a faithful lower-semicontinuous structure).
\field{Configuration} The TV field, and the poset-aware (partially-ordered-label)
field.
\field{Statistic} Number of spurious islands (a high-dimensional anchor all of
whose graph neighbors have strictly lower fitted dimension, away from the true
high-dimensional region).
\field{Acceptance} The poset-aware field has zero spurious islands; the TV field's
island count is reported (it may permit a few; document the gap that motivates the
poset-aware solve).
\field{Safeguards} The cone has a known frontier; islands are counted only away
from the genuinely high-dimensional bulk.

\subsection{T9 --- downstream accuracy (the field is a means, not an end)}
\field{Claim} An LPS using the regularized dimension field attains lower
Truth-RMSE than LPS using independent per-anchor dimension or a fixed global
dimension.
\field{DGP} The stratum-mixture generator carrying a known response surface
$f$ (continuous and binary variants), so probability/continuous Truth-RMSE is
defined.
\field{Configuration} Three LPS arms identical except for the chart-dimension
policy: (i) independent per-anchor, (ii) fixed global, (iii) the TV field; the
field's $\lambda$ and scale selected by held-out Truth-RMSE.
\field{Statistic} Truth-RMSE on held-out points (and, for binary, probability
Truth-RMSE), over $R\ge 30$ replicates, paired.
\field{Acceptance} The TV-field arm is below both baselines by a predeclared
margin; ties or losses on near-manifold (homogeneous-dimension) generators are
acceptable and expected.
\field{Safeguards} Match geometry, response realization, folds, support, kernel,
degree, ridge across arms; vary only the dimension policy. Selecting $\lambda$ by
held-out predictive error (not by an intrinsic dimension score) is mandatory ---
the field exists to help prediction.

\begin{plain}
The tests walk up a ladder: the solver is exact (T1); the right penalty keeps
edges sharp (T2, T4) while cleaning noise (T3); the estimate respects scale (T5),
the right coordinate (T6), the support cap and soft face (T7), and the
geometry's pecking order (T8); and, finally, all of it actually makes the smoother
more accurate (T9). T9 is the one that matters in the end --- the rest exist to
make T9 trustworthy.
\end{plain}

%==============================================================================
\section{Part B --- Experiment specifications (real VALENCIA, exploratory)}
\label{sec:experiments}
%==============================================================================

These run on the France et al.\ VALENCIA 16S collection ($n\approx 13{,}231$
vaginal samples), restricted to \emph{relative abundances of phylotypes} --- each
sample closed to sum one, so it is a point of the simplex and nothing else
(no host covariates, no counts beyond what defines presence). They have no ground
truth; each produces tables, figures, and a decision rule for what to build or
prioritize next. Each follows the template \textsf{Objective / Procedure /
Outputs / Decision}.

\subsection{E1 --- Enumerate the naive strata}\label{sec:E1}
\field{Objective} Map the face-lattice structure the data actually occupy: which
support patterns (naive strata) occur, how heavily, and at what face dimensions.
\field{Procedure}
\begin{enumerate}[leftmargin=1.5em,nosep]
\item Close each sample to the simplex. Fix a presence threshold $\tau$ and define
the support $A_\tau(x)=\{i:x_i>\tau\}$; \emph{sweep} $\tau$ over a grid (e.g.\
$0$, and relative-abundance / read-depth-aware thresholds) because ``present'' is
not crisp at the detection limit.
\item Enumerate the distinct support patterns $\{A\}$; for each record occupancy
$n_A$, size $|A|$, and nominal face dimension $|A|-1$.
\item Refine: group strata by their $L^\infty$ cell (dominant/argmax taxon), and
cross-tabulate against the VALENCIA community-state-type (CST) label.
\end{enumerate}
\field{Outputs} An occupancy-ranked table of strata; the face-dimension
distribution; the count of distinct strata and the occupancy concentration (how
few strata hold, say, $90\%$ of samples) as a function of $\tau$
(Figure~\ref{fig:fig_strata_enum.pdf}); the $L^\infty$-cell and CST cross-tabs.
\begin{decision}
If a small number of strata hold most of the mass, focus all synthetic modeling
and tests on those. If the face-dimension distribution is concentrated at low
dimensions, the non-manifold structure is mild and a global chart dimension may
suffice (\emph{lowers} the priority of the whole field idea). If it is broad,
local, varying dimension is warranted. The $\tau$-sensitivity quantifies how
``soft'' the faces are --- how much the stratum assignment moves as the threshold
moves --- and therefore how much the soft-face machinery (T7) is needed.
\end{decision}

\fig{fig_strata_enum.pdf}{\linewidth}{Schematic of the E1 outputs. (a) Each sample
carries a support pattern --- which phylotypes are present --- equivalently which
face of the simplex it sits on. (b) Occupancy of the naive strata, expected to be
heavy-tailed (a few dominant patterns). (c) The distribution of nominal face
dimension $|A|-1$. (Illustrative values, not computed from the data.)}

\subsection{E2 --- Mass concentration around lower-dimensional faces}\label{sec:E2}
\field{Objective} For each well-occupied stratum, quantify how far \emph{below} its
nominal face dimension the data actually live --- the gap that the local-dimension
field exists to recover.
\field{Procedure} For each stratum $A$ with $n_A\ge n_{\min}$, restrict to its
samples and work in the cube coordinate $\xi$ of the present parts (and, for the
T6 contrast, the log-ratio $u$):
\begin{enumerate}[leftmargin=1.5em,nosep]
\item \emph{Effective rank} $k_A$: the eigenvalue spectrum of the within-stratum
covariance; report the scree and a gap-based $k_A$, alongside the nominal $|A|-1$.
\item \emph{Per-sample local effective dimension}: local PCA at the bandwidth
scale, giving a distribution of effective dimension \emph{within} the stratum.
\item \emph{Near-face mass}: for each lower face $B\subset A$, the fraction of
stratum-$A$ samples within $\epsilon$ of $B$, i.e.\ with $\min_{i\in A\setminus
B}\xi_i<\epsilon$; sweep $\epsilon$.
\item \emph{Min present-abundance law}: the distribution of $\min_{i\in A}x_i$
(how often a present taxon hugs the detection limit), and the within-stratum
anisotropy (largest/smallest eigenvalue $=$ a tube-thickness estimate $\rho_A$).
\end{enumerate}
\field{Outputs} Per stratum: the (nominal $|A|-1$) vs.\ (effective $k_A$) gap
(Figure~\ref{fig:fig_eff_dim_vs_nominal.pdf}); near-face mass curves and
min-abundance histograms (Figure~\ref{fig:fig_mass_conc.pdf}); the tube-thickness
estimates $\rho_A$.
\begin{decision}
The size of the (nominal $-$ effective) gap, per stratum, is the central number.
Large gaps mean the data genuinely concentrate below their face dimension, so the
effective-dimension reduction is high-value and the field should be built and
tested in earnest; the measured $k_A$ and $\rho_A$ become the parameters of the
synthetic generators (E3, Section~\ref{sec:synthetic}). Gaps near zero mean the
faces are ``full'' and the support cap (T7) already suffices --- deprioritize the
within-face reduction and keep only the soft-face and the cap.
\end{decision}

\fig{fig_eff_dim_vs_nominal.pdf}{0.62\linewidth}{Schematic of the headline E2
statistic. For each stratum, the support gives a nominal dimension (the cap, on the
diagonal); the data fill fewer effective directions (below the diagonal). The
vertical gaps are exactly what local-dimension regularization is meant to recover.
(Illustrative.)}

\fig{fig_mass_conc.pdf}{\linewidth}{Schematic of the per-stratum E2 diagnostics: a
naive face whose mass hugs a lower face (a tube); the distribution of the smallest
present abundance (mass near a sub-face when it piles up at the detection limit);
and the within-stratum effective-dimension distribution sitting below the nominal
value. (Illustrative.)}

\subsection{E3 --- Remaining structure for synthetic modeling}\label{sec:E3}
\field{Objective} Collect the rest of the statistics needed to make the synthetic
generators realistic.
\field{Procedure} Estimate, from the data: the stratum distribution $p(A)$ (E1);
per-stratum location and shape --- mean $\mu_A$ and covariance $\Sigma_A$ in a
log-ratio chart (a logistic-normal fit), from which the top-$k_A$ loading
directions $B_A$ and the residual thickness $\rho_A$ follow (E2); dominant-taxon
($L^\infty$-cell) frequencies and the near-tie frequency (how often the top two
taxa are close --- the soft-cell regime); pairwise presence co-occurrence (the
support-layer structure, shared with the chart-aware association design); and the
same quantities conditional on CST.
\field{Outputs} A single parameter bundle $\{p(A),\mu_A,B_A,\rho_A,k_A,\text{
dominant-taxon and near-tie frequencies, co-occurrence}\}$ that fully specifies
the generators of Section~\ref{sec:synthetic}.
\begin{decision}
Before any test built on the generators is trusted, run a posterior-predictive
fidelity check: simulate from the bundle and compare synthetic vs.\ real marginals
(per-taxon abundance laws, stratum occupancy, dominant-taxon frequencies,
co-occurrence). If the synthetic data does not reproduce the real marginals, the
generators --- and the example set built on them --- are not yet realistic, and the
bundle must be refined before the test outcomes mean anything about VALENCIA.
\end{decision}

\begin{plain}
The experiments are a measuring exercise, not a graded one. We are taking the
real microbiome data apart to answer: how many distinct ``which-taxa-are-present''
patterns are there, how crowded is each, and within each, do the present taxa
really spread out or do they huddle near an edge? Those measurements are the
recipe we then use to cook realistic fake data --- and a good cook tastes the dish
(the fidelity check) before serving it.
\end{plain}

%==============================================================================
\section{Part C --- The VALENCIA-inspired synthetic example set}\label{sec:synthetic}
%==============================================================================

These generators expand the non-manifold example set used in previous LPS tests
(replacing the toy ``edge-glued-to-sheet'' object with simplex-resident,
VALENCIA-shaped objects). Every one carries a \emph{known} local-dimension field,
which is what makes the tests of Section~\ref{sec:tests} possible.

\subsection{The stratum-mixture generator (the realistic workhorse)}
\begin{dgp}[Stratum mixture]\label{dgp:mix}
Fix the parameter bundle from E3. To draw one sample:
\begin{enumerate}[leftmargin=1.5em,nosep]
\item draw a support pattern $A\sim p(A)$ with $|A|=m$ present parts;
\item in a log-ratio chart of the present parts (ILR, or the $L^\infty$
max-reference cube coordinate of the companion design), $\R^{m-1}$, draw
\[
  y \;=\; \mu_A \;+\; B_A\,z \;+\; \rho_A\,\varepsilon,
  \qquad z\sim\mathcal N(0,I_{k_A}),\ \ \varepsilon\sim\mathcal N(0,I_{m-1}),
\]
with $B_A\in\R^{(m-1)\times k_A}$ of orthonormal columns (the face's $k_A$ active
tangent directions) and $\rho_A$ small (the tube thickness). The chart covariance
is $\Sigma_A = B_AB_A^{\!\top}+\rho_A^2 I$, with $k_A$ large eigenvalues and
$m-1-k_A$ small ($\rho_A^2$) ones;
\item map $y$ back to the simplex face $F_A$ (inverse chart), set absent parts to
$0$, and close to sum one.
\end{enumerate}
\textbf{Ground truth.} The naive stratum is $A$; the effective local dimension at
the smoother's bandwidth is $k_A$ (we chose it). At a scale finer than $\rho_A$ the
dimension reads $m-1$, coarser than $\rho_A$ it reads $k_A$ --- the designed
scale step that T5 checks. A response surface for T9 is attached on the latent
factors, e.g.\ continuous $f=g(z)$ or binary $p=\mathrm{expit}(\eta(z))$.
\end{dgp}

\subsection{Single-phenomenon fragments (clean unit tests)}
Each fragment isolates one effect with a known field, for the corresponding test.

\begin{dgp}[Tube]\label{dgp:tube} One stratum, $m=3$, $k_A=1$: a curve inside a
$2$-face, thickness $\rho$ tunable. Feeds T5 (scale) and T6 (coordinate trap; place
the curve heading toward a $1$-face so $\xi$ and $u$ disagree).\end{dgp}

\begin{dgp}[Seam]\label{dgp:seam} Two strata of dimensions $k_1\neq k_2$ sharing a
sub-face, glued so that crossing the shared face the true dimension jumps along a
\emph{known} boundary. Feeds T2 and T4 (the boundary tests); the jump magnitude
$|k_1-k_2|$ is a knob.\end{dgp}

\begin{dgp}[Cone]\label{dgp:cone} A stratum whose effective dimension decreases
continuously to a singular point (a $2$-region pinching toward a vertex), a
faithful lower-semicontinuous structure with a known frontier. Feeds T8 (the
no-island / poset test).\end{dgp}

\begin{dgp}[Near-tie ridge]\label{dgp:tie} Two dominant taxa held near-equal so the
$L^\infty$ argmax is unstable, populating the soft-cell regime. Stresses the
interaction of the dimension field with the soft atlas; an auxiliary generator for
robustness checks.\end{dgp}

\fig{fig_synth_model.pdf}{\linewidth}{The stratum-mixture generator
(DGP~\ref{dgp:mix}). A stratum is drawn from the E1 occupancy; a composition is
sampled on that face with a covariance of controlled effective rank $k_A$ (a known
tube); the output carries its true stratum $A_j$ and true local dimension $k_j$, so
the dataset is non-manifold \emph{and} gradeable. (Schematic.)}

\begin{plain}
We build the fake data the way the real data is built --- pick which microbes are
present, then spread the sample out on that face --- but with one twist we control:
we decide, for each face, how many directions the data really spreads in (the
``effective dimension''), and how thin the leftover directions are (the ``tube'').
Because we set those numbers, we know the answer key. The fragments are stripped-
down versions, each built to make exactly one test fair: a single thin curve, a
single seam where the dimension jumps, a single cone that pinches to a point.
\end{plain}

%==============================================================================
\section{Part D --- Sequencing and the decision map}\label{sec:sequencing}
%==============================================================================

\begin{enumerate}[leftmargin=1.6em]
\item \textbf{E1 first} (cheap). It alone can lower the priority of the whole
program: if the real data's face-dimension structure is mild, a global chart
dimension may be enough and the field is not worth building.
\item \textbf{E2 next} --- the decisive measurement. The (nominal $-$ effective)
gap is the quantity that justifies, or deflates, the local-dimension field. Its
$k_A,\rho_A$ feed the generators.
\item \textbf{E3 + fidelity check.} Assemble the parameter bundle; verify the
generators reproduce the real marginals before trusting any downstream test.
\item \textbf{Build the generators} (DGPs~\ref{dgp:mix}--\ref{dgp:tie}).
\item \textbf{Run the tests}, T1 first (a pure unit gate on the solver), then
T2--T8 (the property tests), and T9 last (the payoff: does the field improve LPS
accuracy on VALENCIA-realistic data?).
\end{enumerate}

The outcomes map to next steps. If E2 shows large gaps \emph{and} T9 shows the
field helps on the realistic generator, promote \texttt{local.auto} plus the
$\ell_1$ field to a routine policy and proceed to a real-data LPS evaluation
(where the field is judged only by held-out predictive accuracy, never by an
intrinsic dimension score). If E2 shows small gaps, keep only the support cap and
soft face and shelve the within-face reduction. If T8 shows the TV field
producing geometrically impossible islands on the cone, invest in the
poset-aware solve; otherwise TV suffices. And throughout, the synthetic example
set built here becomes the permanent regression suite for the dimension field,
just as the toy stratified object was for the earlier LPS tests --- only now
shaped like the data we actually care about.

\begin{decision}
The single go/no-go gate for the whole effort is the conjunction
\emph{(E2 gap is large)} $\wedge$ \emph{(T9 improves Truth-RMSE on the
stratum-mixture generator)}. If either fails, the local-dimension field is not
worth carrying into production for vaginal 16S, and the effort redirects to the
cheaper soft-face-and-cap refinement of the chart-aware design.
\end{decision}

%==============================================================================
\begin{thebibliography}{9}\small
\bibitem{expo} \emph{Local Dimension Regularization in Local Polynomial
  Smoothing} (the expository note),
  \nolinkurl{~/current_projects/geosmooth/dev/notes/foundations/lps_local_dimension_regularization.tex}.
\bibitem{plan} \emph{LPS Correctness and Validation: A Precise Experimental Plan},
  \nolinkurl{~/current_projects/geosmooth/dev/methods/lps/specs/lps_experimental_plan_2026-06-09.tex}.
\bibitem{chartaware} \emph{Chart-Aware Local Association Structure on the
  $L^\infty$ Atlas},
  \nolinkurl{~/current_projects/geosmooth/dev/methods/lcov/specs/chart_aware_local_association_design_2026-06-09.tex}.
\bibitem{france2020} M.~France et al., VALENCIA: a nearest-centroid classification
  method for vaginal microbial communities, \emph{Microbiome} 8:166 (2020).
\bibitem{ravel2011} J.~Ravel et al., Vaginal microbiome of reproductive-age women,
  \emph{PNAS} 108 (2011).
\bibitem{ishikawa} H.~Ishikawa, Exact optimization for Markov random fields with
  convex priors, \emph{IEEE TPAMI} 25 (2003).
\bibitem{fusedlasso} R.~Tibshirani, M.~Saunders, S.~Rosset, J.~Zhu, K.~Knight,
  Sparsity and smoothness via the fused lasso, \emph{JRSS-B} 67 (2005).
\bibitem{levinabickel} E.~Levina, P.~J.~Bickel, Maximum likelihood estimation of
  intrinsic dimension, \emph{NeurIPS} (2004).
\bibitem{ilr} J.~J.~Egozcue et al., Isometric logratio transformations for
  compositional data analysis, \emph{Math.\ Geol.} 35 (2003);
  J.~Aitchison, \emph{The Statistical Analysis of Compositional Data} (1986).
\end{thebibliography}

\end{document}
